\section{The problem and the approach}\label{ms:introduction}

The central limit theorem describes the shape of a long sum. The
Berry--Esseen inequality asks how far that description can be from the
truth at a specified sample size. Even when the summands have the same
distribution, the best universal constant is a delicate question: the
largest errors come from the interaction of skewness with the jumps of a
lattice distribution. A proof of the sharp inequality must retain both
effects without allowing exceptional laws to escape as the sample size
changes.

Let $\mathcal P$ denote the probability laws on $\R$ satisfying
$\E X=0$, $\E X^2=1$, and $\rho(P):=\E|X|^3<\infty$. For independent
$X_1,\ldots,X_n$ with common law $P$, write $S_n=X_1+\cdots+X_n$ and
\begin{align*}
 \Delta_n(P)&=\sup_{H\in\R}
       |\Pp_P(S_n\le H)-\Phi(H/\sqrt n)|,\\
 R_n(P)&=\frac{\sqrt n\,\Delta_n(P)}{\rho(P)},\qquad
 C_n=\sup_{P\in\mathcal P}R_n(P),\qquad C^*=\sup_{n\ge1}C_n.
\end{align*}
Here $\Phi$ and $\phi$ are the standard normal distribution function and
density. The conjecture is
\begin{equation}\label{ms:conjecture}
 C^*=\CE,\qquad
 \CE=\frac{3+\sqrt{10}}{6\sqrt{2\pi}}=0.409732\ldots.
\end{equation}
Throughout, $H$ is a threshold for the \emph{raw sum}; $H/\sqrt n$ is its
standardized version. A later Bernoulli coordinate introduces a third
threshold, and we will keep the three distinct. All terminating decimals
in numerical inequalities mean exact rational numbers.

The normalization costs no generality. If $Z$ has mean $\mu$, variance
$\sigma^2>0$, and finite third absolute central moment $\beta_3$, apply the
definitions to $X=(Z-\mu)/\sigma$. Then $\rho=\beta_3/\sigma^3$ and
\[
 \Pp\left(\frac{\sum_{i=1}^n Z_i-n\mu}{\sigma\sqrt n}\le x\right)
 =\Pp(S_n\le\sqrt n\,x).
\]
Thus $C_n\le\CE$ is precisely the usual sharp inequality at index $n$.

\subsection{From Berry and Esseen to the sharp-constant problem}

Berry~\cite{Berry1941} and Esseen~\cite{Esseen1942,Esseen1945}
established quantitative normal approximation of order $n^{-1/2}$ under
a finite third absolute moment. Their results initiated the search for
the smallest constant in a bound that is uniform over the summand law.
Esseen's later moment analysis~\cite{Esseen1956} identified the lower
constant in~\eqref{ms:conjecture}. The conjecture that this lower constant
is also the universal upper constant is associated with Zolotarev's work
of 1966--1967~\cite{Zolotarev1966,Zolotarev1967}.

The two-point example explains the value of the constant. Put
\[
 p_*=(4-\sqrt{10})/2,\qquad q_*=1-p_*,\qquad
 X=(B-p_*)/\sqrt{p_*q_*},\quad B\sim\operatorname{Bernoulli}(p_*).
\]
For this fixed law, $\rho=(p_*^2+q_*^2)/\sqrt{p_*q_*}$ and the
classical lattice asymptotic gives
\begin{equation}\label{ms:bernoulli-limit}
 \lim_{n\to\infty}R_n(P)
 =\frac{3+q_*-p_*}{6\sqrt{2\pi}(p_*^2+q_*^2)}=\CE.
\end{equation}
Indeed $p_*^2+q_*^2=10-3\sqrt{10}$ and
$(3+q_*-p_*)/(p_*^2+q_*^2)=3+\sqrt{10}$.
Consequently every constant smaller than $\CE$ fails for a sufficiently
large sum of this one fixed Bernoulli law. This proves $C^*\ge\CE$;
it supplies no upper bound for other laws. Schulz's dissertation
\cite{Schulz2016} proves the sharp upper bound for the entire binomial
family at every sample size. That finite-binomial theorem and the
fixed-law lower asymptotic have different roles in the argument below.

Upper estimates developed through increasingly precise smoothing and
characteristic-function inequalities. Zolotarev's work and van Beek's
Fourier analysis~\cite{Zolotarev1967,vanBeek1972} are early milestones;
Prawitz's one-sided smoothing kernel~\cite{Prawitz1972} remains useful in
modern computations. More recent improvements include Shevtsova's bound
$0.7056$~\cite{Shevtsova2007}, the bound $0.5129$ of Korolev and
Shevtsova~\cite{KorolevShevtsova2010}, Tyurin's bound $0.4785$ obtained
using zero-bias methods~\cite{Tyurin2009}, and Shevtsova's bounds
$0.4748$ and $0.4690$~\cite{Shevtsova2011,Shevtsova2013}.
These estimates also give useful information before the absolute constant
is extracted. For example, the structural inequalities
\begin{equation}\label{ms:structural-classical}
 \sqrt n\,\Delta_n(P)\le
 \min\{0.4690\rho,\ 0.3322(\rho+0.429),\
                          0.3031(\rho+0.646)\}
\end{equation}
hold for every $P\in\mathcal P$ and $n\ge1$~\cite{Shevtsova2013}.
If $R_n(P)$ is near $\CE$, they bound $\rho$ from above. This supplies the
first quantitative moment cap for the support calculation.

\subsection{Recent manuscripts, repositories, and formalization}

Several recent projects combine analytic inequalities with extensive
computation. Xiao and Li report a global iid upper bound of $0.4395$;
Lin reports $0.423$ using a refinement of their analytic framework
\cite{xiao2026berryesseen,lin2026berryesseen}. Both supply Lean
developments and exact finite certificates. Their use of
\texttt{native\_decide} places the relevant compiler and evaluation
mechanism within the trusted computation. We have inspected the stated
scope and verification descriptions, but have not independently replayed
these developments; their numerical conclusions are not premises of this
paper. Tao and contributors' \emph{Optimization Constants in Mathematics},
especially Problem~19a, records this work and its verification status
\cite{tao2026berryesseenregistry}. The registry is a useful guide to the
literature, rather than an additional proof of a listed bound.

He and Cheng's preprint~\cite{hecheng2026eventually} reports the sharp
inequality for every
\[
 n\ge 2\lceil\exp(10^{17})\rceil,
\]
with a threshold independent of the summand law. This uniform quantifier
is essential. A theorem for each fixed law, with a law-dependent
threshold, would not exclude a sequence of counterexamples whose laws
change with $n$. Their eventual theorem also leaves a finite range to be
addressed; it is not an all-index proof. We cite version~3 explicitly,
since the title and argument changed between versions.

Bell's collection considers independent summands without the iid
restriction. It reports a universal upper bound below $0.44988$, finite
array results, and sharp bounds in a small-Lyapunov regime
\cite{BellIndependent2026,BellFinite2026,BellSmall2026}.
The collection's extremizer and structural papers
\cite{BellExtremizers2026,BellStructural2026} distinguish the distribution
of a sum from the third-moment cost of an independent factorization of
that distribution. They also distinguish an extremal value from a
classification or uniqueness theorem. These distinctions matter here:
a necessary contact condition does not establish that a relaxed optimizer
is an iid convolution. In the independent setting, zero padding can also
change the number of summands without changing the normalized sum, so
sample-size assertions must be compared with care. The cited papers are
from collection version~0.7.0; the Primal Space project page describes an
earlier version of the same collection, not a separate result.\footnote{Project overview: \url{https://www.primalspace.me/projects/independent-berry-esseen-bounds/}.} We have not
replayed its computational certificates.

Formalization offers another way to expose the exact hypotheses of a
result. Yan's \emph{ProbabilityApproximation} formalizes classical
nonuniform scalar and multivariate normal-approximation theorems in Lean
\cite{Yan2026ProbabilityApproximation}; it is not a formal proof of the
sharp iid constant. The \emph{BerryEsseenLean} project accompanying He
and Cheng has eventual-sharpness endpoints with explicitly listed external
mathematical inputs~\cite{BerryEsseenLean2026}. Its inspected September
snapshot accompanies an earlier manuscript version. A kernel-checked
implication with such inputs is valuable, but it does not discharge those
inputs. We have not rebuilt either formalization.

\subsection{What this manuscript proves and how to read it}

Our aim is to reduce the analytic and computational cost of a possible
all-index proof, while making the remaining obligations precise. An
earlier working draft of this project is available on Hexagon
\cite{rodziewiczkolodziejek2026momentdefects}. We use it as source material,
not as independent confirmation of its conclusions. The present draft
does \emph{not} claim~\eqref{ms:conjecture}.

The central simplification is to examine the first index at which a
sharp bound could fail. If a base through $41$ is available, that index
is at least $42$, and the immediately preceding global constant is at
most $\CE$. The resulting contact inequality has no positive predecessor
loss. A support-free smoothing estimate also bounds the current score; the least-index choice supplies the
bounds at both deleted indices. This is the role of
Section~\ref{ms:smoothing}; it does not by itself establish the finite
base or the local exclusion theorem.

The proof then has a concrete geometric form. A maximizing law must
satisfy a contact inequality at every support point. That inequality
confines the law to a bounded interval. Splitting the law by sign gives
a two-point skeleton with conditional mean-zero noise. A nonnegative
moment defect measures how much probability can lie away from the
skeleton; one joint inequality controls both the noise variance and the
inward displacement of the two support clusters. Finally, conditioning
on the number of positive labels expresses the original iid sum exactly
as a binomial mixture. Local comparisons can then be certified on closed
parameter cells, provided that all these cells refer to the same
underlying law.

Sections~\ref{ms:support}, \ref{ms:resources}, and \ref{ms:curvature}
give three reductions in this program. The support argument replaces
polynomial positivity meshes by two rational boundary inequalities and
one monotonicity calculation. Keeping the score excess in the joint
defect bound makes several comparison sectors empty. An exact
one-dimensional maximization replaces a curvature grid by at most three
candidates. Each statement carries its hypotheses, including any
previously acquired moment or phase bounds.

Section~\ref{ms:certificates} states the conditional global implication
and separates analytic inputs from finite certificate conditions.
The appendices contain the longer retained derivations, with their
original numerical scopes made explicit. Some belong to a supplied
proof candidate whose full independent verification remains open; they
are marked accordingly. This distinction is particularly important when
an argument proved for $n\ge47$ is proposed for use at $n\ge42$.
The source correspondence and citation audit accompany the manuscript
as review files, so that a reader can trace every retained dependency.
